% ============================================================
%  Chapter 1 — Introduction
% ============================================================
\chapter{Introduction}
\label{ch:introduction}

\noindent The International Data Corporation (IDC)~\cite{idc:2021} estimates that by 2025, \ac{IoT} devices---encompassing machines, sensors, and cameras---will collectively generate 79.4~zettabytes of data.
Although \acp{SD} can increase consumer well-being~\cite{lee:lee:2020}, they also present significant challenges.
Sensors invade the consumer privacy sphere, capturing data stored in silos that attract cybercriminals~\cite{weber:2010}, while \acp{OEM} often exploit these data for their own benefit~\cite{zuboff:2019}, leaving consumers without tangible rewards~\cite{spiekermann:korunovska:2017}.
The framework governing the collection and use of personal data involving consumers and \acp{OEM} to create economic value, known as the \ac{PDE}~\cite{world:economic:forum:2011}, is unbalanced~\cite{ctrl:shift:2014}.
It allows \acp{OEM} to monopolize data-flow profits while consumers bear the system's externalities, and is therefore associated with a loss of control over personal information~\cite{abiteboul:andre:abiteboul:2015}.

To address this imbalance, industry and academia are exploring blockchain-based \ac{IoT} solutions for people-centered architectures~\cite{samaniego:deters:2017}.
Blockchain's strong cryptographic foundation, immutable and tamper-proof nature, decentralization, and smart-contract support allow for critical mechanisms that, when integrated with user-centered architectures, elevate consumer status in the \ac{PDE}~\cite{zyskind:nathan:pentland:2015}.
By leveraging blockchain \ac{DLT}, it is possible to establish a network of trustless peers that collectively agree on an immutable, auditable, and cryptographically secure shared version of the truth~\cite{zheng:xie:dai:chen:wang:2018}.
This network offers support for decentralized identifiers that link users with their data without the need for third parties~\cite{muhle:suhr:zanker:meinel:2018}.
When combined with Personal Data Stores, frameworks that allow users to collect, store, manage, and share their data according to their preferences~\cite{deangelis:pieri:2019}, this identity mechanism, referred to as \ac{SSI}~\cite{preukschat:reed:2021}, gives consumers sovereign control over their data.
When used within a user-centered framework, this mechanism transforms consumers from passive data subjects to active participants in the \ac{PDE}, underscoring the critical role of maintaining data rights---access, usage, and sharing---firmly in the hands of consumers~\cite{janssen:brous:estevez:barbosa:2018}.

The \ac{DT} concept originated at NASA and became publicly recognized in 2003~\cite{grieves:2014}.
It refers to a virtual representation or digital counterpart of a physical object, system, or process~\cite{minerva:lee:crespi:2020}.
Since its introduction, the concept has evolved and matured, reaching a significant growth stage by 2014~\cite{grieves:vickers:2017}.
Transforma Insights~\cite{transforma:insights:2021} forecasts a market size of 24.1~billion devices and a reach of \$1.5~trillion, with the consumer sector accounting for 65\% of all connections.
\ac{DT} reframes much of the early efforts to develop connected-device solutions, initially based on telemetry and \ac{M2M} communications and later enhanced with \ac{IoT} technologies~\cite{madakam:ramaswamy:tripathi:2015}.
This evolution has produced a more structured, multifaceted development framework supporting services and \ac{AI}~\cite{minerva:lee:crespi:2020}.
For instance, a smartwatch with a \ac{DT} extends beyond basic data collection: it continuously gathers and analyzes detailed data to update and refine its virtual model, creating a comprehensive understanding of the user's behaviors, activities, and physiological states.

As \ac{DT} crosses into the consumer realm, it is reasonable to surmise that some \acp{OEM} may exploit the concept of product personalization~\cite{thiesse:2007} or other features with disproportionately lower impact for consumers to gain further consumer insights with \ac{DT}-based solutions.
Therefore, it is vital to develop mechanisms that enable consumers to control their devices' \acp{DT}~\cite{kopponen:hahto:kettunen:etal:2022}.
Tesla is a paradigmatic example of the urgent need for such mechanisms: its trillion-dollar valuation~\cite{tesla:marketcap:2024} is supported by its ability to leverage \acp{DT} to train its Full Self-Driving feature~\cite{karpathy:matroid:2020,tesla:ai:day:2021}.

%--------------------------------------------------------------
\section{Problem Statement}
\label{sec:problem}
%--------------------------------------------------------------

\noindent The economic fairness challenge in controlling \acp{DT} arises from the current cloud-based approach, where \acp{OEM} retain \emph{de facto} control over the \ac{DT} and the associated consumer-generated data~\cite{hummel:braun:dabrock:2021}.
As consumers become more aware of the growing value of data, they are likely to demand solutions that restore control and ensure equitable participation.

Furthermore, proponents of edge computing, federated learning, and privacy-preservation methods such as homomorphic encryption emphasize the benefits of storing and processing data at the edge as a mechanism to improve the privacy and security of \ac{AI} training data within \ac{IoT} contexts~\cite{idc:ai:2024,shi:cao:zhang:li:xu:2016,mcmahan:moore:ramage:hampson:arcas:2017,yang:liu:chen:tian:2019,li:sahu:talwalkar:smith:2020}.
However, these works often overlook the critical step of ensuring that data are kept at the edge in a self-sovereign manner under consumer control.

This dissertation addresses these challenges by presenting the \ac{C2DTA}.
The main objective of \ac{C2DTA} is to empower consumers in the \ac{PDE} by shifting the \ac{DT} from the cloud under \ac{OEM} control to the edge under consumer control~\cite{spiekermann:korunovska:boenisch:2019}.
This shift makes consumers \emph{de facto} controllers---meaning they have actual control in practice~\cite{finck:2018}---of their \ac{SD} data by giving them effective control over the \ac{DT} and its associated data~\cite{de:hert:papakonstantinou:2016}.

\ac{C2DTA} addresses two key challenges:
\begin{enumerate}
    \item \textbf{Economic fairness through the control of \acp{DT}.} The current cloud-based approach allows \acp{OEM} to retain \emph{de facto} control over the \ac{DT} and associated consumer-generated data. \ac{C2DTA} restores consumer sovereignty by hosting the \ac{DT} at the edge.
    \item \textbf{Bridging the edge-computing gap.} \ac{C2DTA} provides the practical foundation necessary for self-sovereign, edge-based data management, closing the gap between proposals for edge computing, federated learning, and privacy-preserving technologies and their actionable implementation.
\end{enumerate}

%--------------------------------------------------------------
\section{Research Objectives}
\label{sec:objectives}
%--------------------------------------------------------------

\noindent Building on the problem statement, this dissertation pursues the following research objectives:

\begin{description}
    \item[RO1.] \textbf{Self-sovereign \ac{DT} ecosystem architecture.} Propose an architecture that enables stakeholders of the digital-twinned \ac{SD} ecosystem to securely manage \acp{DT} and leverage their data. In this architecture, those who own devices control the associated \acp{DT} at the edge, and stakeholders transact in a decentralized, self-sovereign, and secure manner.

    \item[RO2.] \textbf{\ac{DT}, \ac{SSI}, and blockchain integration taxonomy.} Propose a structured framework for consistently evaluating conceptual and practical implementations, enabling clear classification, identifying research gaps, and guiding future development in decentralized systems.

    \item[RO3.] \textbf{Hyperledger Fabric, Indy, and Aries integration.} Propose an architecture that integrates Hyperledger Fabric for managing \acp{SD} and datasets with Hyperledger Aries for overseeing stakeholder identifiers, \acp{VC}, and secure communication protocols through ACA-Py. This integration encompasses eight detailed use cases across the full \ac{SD} lifecycle---manufacturing, sale, ownership transfer, twinning, and untwinning.

    \item[RO4.] \textbf{Future innovation for consumer empowerment.} Introduce concepts emerging from the architecture, including ``data-only \acp{DT}'' and the aggregation of several consumer-controlled \acp{DT} into the Consumer-Controlled Personal Digital Twin (\ac{C2PDT}).
\end{description}

%--------------------------------------------------------------
\section{Research Methodology}
\label{sec:methodology}
%--------------------------------------------------------------

\noindent This dissertation adopts \ac{DSR} as its overarching research methodology.
\ac{DSR} is a problem-solving paradigm rooted in engineering and the sciences of the artificial~\cite{simon:1996} that seeks to extend the boundaries of human and organizational capabilities by creating new and innovative artifacts~\cite{hevner:march:park:ram:2004}.

Following the \ac{DSR} methodology process model proposed by Peffers et al.~\cite{peffers:tuunanen:rothenberger:chatterjee:2007}, the research proceeds through six activities:

\begin{enumerate}
    \item \textbf{Problem identification and motivation.} The imbalanced \ac{PDE}, where \acp{OEM} monopolize consumer data generated by \acp{SD}, is identified as the core problem (Chapter~\ref{ch:introduction}).
    \item \textbf{Objectives of a solution.} The research objectives (RO1--RO4) define what the artifact must achieve (Section~\ref{sec:objectives}).
    \item \textbf{Design and development.} The \ac{C2DTA} architecture is designed and implemented as the research artifact, integrating \ac{DT}, \ac{SSI}, and blockchain technologies (Chapter~\ref{ch:architecture}).
    \item \textbf{Demonstration.} The architecture is demonstrated through a business scenario that tracks the lifecycle of a smartwatch equipped with a heartbeat sensor, from manufacture to purchase, twinning, and resale, encompassing eight use cases (Chapter~\ref{ch:results}).
    \item \textbf{Evaluation.} Feasibility testing confirms that \ac{C2DTA} effectively empowers consumers to manage their \ac{SD} \acp{DT} and associated data. Performance benchmarks assess the impact of the decentralized infrastructure (Chapter~\ref{ch:results}).
    \item \textbf{Communication.} The results are communicated through this dissertation and the accompanying publication~\cite{pinto:silva:moro:aquino:2025}.
\end{enumerate}

The artifact type is an \emph{instantiation}, as defined by March and Smith~\cite{march:smith:1995}: a working system that demonstrates that constructs, models, and methods can be implemented, enabling concrete assessment of the artifact's suitability to its intended purpose.

A problem-centered approach~\cite{peffers:tuunanen:rothenberger:chatterjee:2007} is adopted, as the research originates from the observation of the imbalanced \ac{PDE} and the need to empower consumers with control over their \acp{SD}' \acp{DT}.

To ground the state of the art, a \ac{SLR} is conducted following the PRISMA guidelines~\cite{moher:liberati:tetzlaff:etal:2010} and the methodology proposed by Kitchenham and Charters~\cite{kitchenham:charters:2007}.
The \ac{SLR} methodology, detailed in Chapter~\ref{ch:state-of-art}, systematically identifies, evaluates, and synthesizes existing research on the integration of \acp{DT}, \ac{SSI}, and blockchain.

%--------------------------------------------------------------
\section{Contributions}
\label{sec:contributions}
%--------------------------------------------------------------

\noindent This dissertation makes four main contributions:

\begin{enumerate}
    \item \textbf{Self-sovereign \ac{DT} ecosystem architecture.} An architecture that enables stakeholders to securely manage \acp{DT} and leverage their data. Device owners control the associated \acp{DT} at the edge, and stakeholders transact in a decentralized, self-sovereign, and secure manner. The architecture further explores its impact on the \ac{PDE} and the business strategies to support it.

    \item \textbf{\ac{DT}, \ac{SSI}, and blockchain integration taxonomy.} A structured framework for consistently evaluating conceptual and practical implementations of \ac{DT}--\ac{SSI}--blockchain systems. The taxonomy classifies systems along three dimensions: \ac{DT} implementation level (simplified vs.\ comprehensive), \ac{SSI} implementation level (partial vs.\ full), and blockchain ecosystem implementation level (absent vs.\ partial vs.\ integrated).

    \item \textbf{Hyperledger Fabric, Indy, and Aries with ACA-Py integration.} An architecture that integrates Hyperledger Fabric for managing \acp{SD} and datasets with Hyperledger Aries for stakeholder identifiers, \acp{VC}, and secure communication protocols. The integration encompasses eight use cases across 88 steps, including automation mechanisms to safeguard ecosystem integrity.

    \item \textbf{Future innovation for consumer empowerment.} The introduction of ``data-only \acp{DT}'' and the aggregation of several consumer-controlled \acp{DT} into the \ac{C2PDT}, positioning consumer-controlled Personal Digital Twins as pivotal in empowering individuals in the \ac{AI} era.
\end{enumerate}

%--------------------------------------------------------------
\section{Document Structure}
\label{sec:structure}
%--------------------------------------------------------------

\noindent This dissertation is organized as follows:

\begin{description}
    \item[Chapter~\ref{ch:state-of-art} --- State of the Art.] Presents the \ac{SLR} on the convergence of \ac{DT}, \ac{SSI}, and blockchain. Covers \ac{DT} concepts, \ac{SSI} standards, blockchain for \ac{IoT}, the \ac{PDE}, and data sovereignty. Analyzes related work and identifies research gaps.

    \item[Chapter~\ref{ch:architecture} --- \ac{C2DTA} Architecture.] Introduces the \ac{C2DTA} design, detailing the integrated technologies---Eclipse Ditto, Hyperledger Fabric, Hyperledger Indy, ACA-Py, \ac{IPFS}, \ac{MQTT}, and \ac{DIDComm}---and the eight use cases that cover the full \ac{SD} lifecycle.

    \item[Chapter~\ref{ch:results} --- Results and Discussion.] Presents the implementation, the smartwatch case study, feasibility testing, and performance benchmarks. Discusses the architecture's impact on the \ac{PDE} and business strategies.

    \item[Chapter~\ref{ch:conclusions} --- Conclusions and Future Work.] Summarizes the findings, discusses limitations, and outlines future research directions, including the \ac{C2PDT} concept and its role in the \ac{AI} era.
\end{description}
